\section*{第 \ref{ch:b3:topology} 章 --- 一般拓扑}

\begin{solution}{exo:b3:topology:1}
(a) 集合 $\{a,b\}$ 上共有四种拓扑：平凡拓扑
$\{\varnothing, X\}$；离散拓扑；两个\emph{谢尔宾斯基}
拓扑 $\{\varnothing, \{a\}, X\}$ 和 $\{\varnothing,
\{b\}, X\}$。后两种彼此同胚（交换 $a, b$）；由此共有
三类。连通：除离散拓扑外均连通（只有
离散拓扑含有真非空的闭开集）。
豪斯多夫：仅离散拓扑满足（其余拓扑中，$b$ 的唯一
邻域为 $X$，或 $a$ 的唯一邻域）。

(b) $\Q$：内部为 $\varnothing$（每个区间都含有
无理数），闭包为 $\R$（稠密），边界为 $\R$。$[0,1)
\cup \{2\}$：内部为 $(0,1)$，闭包为 $[0,1]\cup\{2\}$，
边界为 $\{0, 1, 2\}$。$\{1/n\}$：内部为 $\varnothing$，
闭包为 $\{1/n\}\cup\{0\}$，边界即为闭包本身。
\end{solution}

\begin{solution}{exo:b3:topology:2}
(a) 每个开集 $V \subseteq Y$ 都是拓扑基中
集合的并 $\bigcup B_j$，于是 $f^{-1}(V) = \bigcup f^{-1}(B_j)$：若后者皆为开集，
则前者亦开；反之显然。

(b) 集合 $f^{-1}\bigl(\intoo{\frac12}{\frac32}\bigr) = [0, \infty)$
不是开集，故 $f$ 不连续；每点的右连续性
显然（$f$ 在每点右侧局部为常数）。
形如 $[a, b)$ 的集合满足拓扑基的判别准则
（见\cref{def:b3:topology:basis}）：$[a,b)\cap[c,d) = [\max(a,c),
\min(b,d))$。索根弗雷拓扑比通常拓扑更细
（因为 $(a, b) = \bigcup_{a < c < b}[c, b)$）；严格地说：
$[0,1)$ 是索根弗雷开集，而非通常开集。关于
索根弗雷直线在 $x$ 处的连续性：$x$ 的索根弗雷邻域
包含某基集 $[a, b) \ni x$，从而包含 $[x, b)$ ---
于是连续性条件即：对每个 $\varepsilon > 0$
存在 $\delta > 0$，使得 $\abs{f(t) - f(x)} < \varepsilon$
只要 $x \leq t < x + \delta$ 即成立。这恰是
$x$ 处的右连续性。
\end{solution}

\begin{solution}{exo:b3:topology:3}
两个非空开集的补集都是有限的，因此它们的
交集的补集也有限：交集非空（该集合
无限） --- 因此任意两点都没有不交的邻域：不是
豪斯多夫的，从而不可度量化
（见\cref{def:b3:topology:hausdorff}）。设 $(x_n)$ 为单射，
$x$ 为任意点：$x$ 的一个邻域包含某个余有限开集
$U$。$X\setminus U$ 的有限个点至
多只对应有限个下标（由单射性），因而数列的
某个尾部落在 $U$ 中，即 $x_n \to x$；对\emph{每一个} $x$ 均如此。
唯一性的证明需要两个不交的邻域来分离
两个假定的极限——这恰恰在此不成立。
\end{solution}

\begin{solution}{exo:b3:topology:4}
(a) 连续性由构造即得
（见\cref{def:b3:topology:constructions}）。开性：设开集 $W
\subseteq X \times Y$ 是矩形 $U_i\times V_i$ 的并，且
$\pi_X(W) = \bigcup U_i$（非空矩形投影到各自的
因子上），故为开。投影不是闭映射：$H = \{(x,y) : xy = 1\}$ 是
$\R^2$ 中的闭集（它是 $\{1\}$ 在连续乘法下的原像），
但 $\pi_X(H) = \R\setminus\{0\}$ 不是闭集。

(b) （$\Rightarrow$）投影是连续的。（$\Leftarrow$）
设 $x^{(k)} \to x$ 按分量收敛，$W = \prod U_n$ 为 $x$ 的一个基
邻域，其中 $U_n = X_n$ 除 $n \in F$ 外成立，
且该例外集有限。对每个 $n \in F$，$x_n^{(k)} \in U_n$ 当 $k \geq
K_n$；当 $k \geq \max_F K_n$ 时 $x^{(k)} \in W$。公式 $d$
确定一个度量（每个求和项的上界为一，并需
验证 $\min(1, d_n)$ 为度量）；其球 $B_d(x,
2^{-N})$ 包含基箱 $\{y : d_n(x_n,y_n) < 2^{-N},\
n \leq N'\}$（对适当的 $N'$，尾部 $\sum_{n > N'} 2^{-n}$
很小），反之每个基箱都包含某个 $d$ 球：
两个拓扑在每点的邻域族相同。

(c) 在乘积拓扑中，由 (b) 知 $x^{(k)} \to 0$。箱式
开集 $U = \prod_n \intoo{-\tfrac1n}{\tfrac1n}$ 含 $0$，但
$x^{(k)} \notin U$（对每个 $k$，第 $n$ 个分量 $1/k
\geq 1/n$ 当 $n \geq k$ 时成立）：无尾部落入 $U$。箱拓扑
严格地更细，并且就收敛而言
已不属乘积型拓扑。
\end{solution}

\begin{solution}{exo:b3:topology:5}
$\varphi \colon \R \to S^1$，$x \mapsto \eu^{2\iu\pi x}$，是
连续、满射，并且在模 $\Z$ 的等价类上取常值：它
诱导出连续双射 $\bar\varphi \colon \R/\Z \to
S^1$（见\cref{prop:b3:topology:universal}(b)）。投影
$\pi$ 是\emph{开映射}：对开集 $U \subseteq \R$，
$\pi^{-1}(\pi(U)) = \bigcup_{n}(U + n)$ 是开集，故 $\pi(U)$ 也
是开的。豪斯多夫性：设 $\bar x \neq \bar y$；取
代表元，使 $\delta = \min_{n}\abs{x - y - n} > 0$；
半径为 $\delta/3$、中心分别为 $x$ 与
$y$ 的区间的像是开集（因 $\pi$ 为开映射），包含 $\bar x, \bar y$，
且不交（若两原像点相距 $< 2\delta/3$
在模 $\Z$ 意义下）。紧性：$\R/\Z = \pi([0,1])$，是
紧空间的连续像（见\cref{thm:b3:topology:compactprops}(2)）。现
$\bar\varphi$ 是从紧空间到
豪斯多夫空间 $S^1$ 的连续双射，从而是同胚
（见\cref{cor:b3:topology:compacthomeo}）。

对于 (iii)：复合映射 $[0,1] \hookrightarrow \R
\xrightarrow{\pi} \R/\Z$ 连续、满射，并且
恰把 $0 \sim 1$ 粘合起来：它诱导出连续双射
$[0,1]/(0\sim1) \to \R/\Z$，定义域为紧空间（即
$[0,1]$ 在商投影下的连续像），到豪斯多夫
空间的连续双射，因而为同胚。合成后，三个空间
彼此同胚。
\end{solution}

\begin{solution}{exo:b3:topology:6}
(a) $K_1 \cup \dots \cup K_m$ 的开覆盖限制到
每个 $K_j$ 上仍为覆盖：每一块取有限个开集即可。
交集 $\bigcap K_i$ 是紧集 $K_{i_0}$ 中的闭集，
（紧集在所在的度量空间中是闭集），因此
故为紧集。

(b) $x \mapsto d(x, F)$ 连续（$1$-利普希茨），且
在 $K$ 上严格为正（$d(x, F) = 0$ 意味着 $x \in \bar F =
F$）；在紧集 $K$ 上它取到最小值 $m > 0$，即 $d(K, F)
\geq m$。若去掉紧性，结论不成立，例如：$F_1 = \{n : n \geq
2\}$ 与 $F_2 = \{n + \frac1n : n \geq 2\}$ 为闭集，
二者不交，且距离为 $\inf_n \frac1n = 0$。

(c) 若 $X$ 紧，则每个连续函数 $f$ 有界
（由\cref{thm:b3:topology:compactprops}(2)）。反之，若 $X$
不紧，取 $(x_n)$ 使其无收敛子列
（见\cref{thm:b3:topology:metriccompact}）；取子列，
不妨设 $x_n$ 两两不同，且
$X$ 中没有点是该数列的聚点，于是
\[
r_n = \min\Bigl(2^{-n},\ \tfrac13\,d\bigl(x_n, \{x_m : m \neq
n\}\bigr)\Bigr) > 0,
\]
球 $B(x_n, r_n)$ 两两不交（若第 $n$ 个与第 $m$ 个球有公共点，
则 $d(x_n, x_m) < r_n +
r_m \leq \frac23 d(x_n, x_m)$）。定义
\[
f(x) = \sum_{n} n\,\max\Bigl(0,\ 1 - \frac{d(x,
x_n)}{r_n}\Bigr):
\]
对每个 $x$，至多有一个加项非零，并且 $f(x_n) = n$。
在 $x$ 处验证连续性：设 $y_k \to x$，每个值 $f(y_k)$
或为 $0$，或为单个凸起值 $\mathrm{bump}_{n_k}(y_k)$。
若某个指标 $n$ 出现无穷多次，则沿该子列
有 $f(y_k) = \mathrm{bump}_n(y_k) \to
\mathrm{bump}_n(x)$，它等于 $f(x)$（若 $x \in B(x_n,
r_n)$，则整个尾部都落在此开球内，不再出现其他
指标；若 $x \notin B(x_n, r_n)$，则
$\mathrm{bump}_n(x) = 0 = f(x)$，因 $x$
属于第 $n$ 个球的闭包，且不在其他任何开球中）。若 $n_k \to \infty$：
$d(y_k, x_{n_k}) < r_{n_k} \leq 2^{-n_k} \to 0$，于是 $x_{n_k}
\to x$，使 $x$ 成为 $(x_n)$ 的聚点 --- 已排除；因此
此情形仅涉及有限个 $k$。总之在每种情形下
$f(y_k) \to f(x)$：$f$ 连续，且无界。
\end{solution}

\begin{solution}{exo:b3:topology:7}
假设不存在这样的 $\delta$：对每个 $n$，存在集合 $A_n$，
其直径 $< 1/n$，且不包含于任何单个 $U_i$；取 $x_n \in
A_n$。由紧性（见\cref{thm:b3:topology:metriccompact}），
存在子列 $x_{n_k} \to x$；取 $i$ 及 $r > 0$，使 $B(x,
r) \subseteq U_i$。当 $k$ 充分大时：$d(x_{n_k}, x) < r/2$ 且
$1/n_k < r/2$，故 $A_{n_k} \subseteq B(x_{n_k}, 1/n_k)
\subseteq B(x, r) \subseteq U_i$ --- 矛盾。海涅：
给定 $\varepsilon$，$Y$ 可由球 $B(y, \varepsilon/2)$ 所覆盖；
其原像构成 $X$ 的开覆盖；取 $\delta$ 为
勒贝格数：若 $d(x, x') < \delta$，则点对 $\{x, x'\}$
直径 $< \delta$，落在某一原像中，且 $d(f(x),
f(x')) < \varepsilon$。
\end{solution}

\begin{solution}{exo:b3:topology:8}
(a) $GL_n(\R) = \det^{-1}(\R\setminus\{0\})$ 为开集（$\det$ 为
多项式，因而连续）。不连通：$\det$ 将其映满
$\R\setminus\{0\}$，而连通空间的连续像仍连通，
（见\cref{thm:b3:topology:connectedbasics}(2)）；
$\R\setminus\{0\}$ 不是区间。$GL_n(\C)$：设 $A, B$
可逆，$p(\lambda) = \det\bigl((1-\lambda)A + \lambda
B\bigr)$ 为 $\lambda$ 的多项式，且 $p(0), p(1) \neq 0$，
故不恒为零：它在 $\C$ 中只有有限个根。
去掉有限个点的复平面是道路连通的（绕开
这些点即可），故存在道路 $\lambda(t)$，
由 $0$ 到 $1$，使 $p(\lambda(t)) \neq 0$：$t \mapsto
(1-\lambda(t))A + \lambda(t)B$ 为 $GL_n(\C)$ 中的道路。

(b) $\det(A + \varepsilon I) = \chi_A(-\varepsilon)\cdot(\pm1)$
至多对 $n$ 个 $\varepsilon$ 取零值：
可逆矩阵 $A + \varepsilon I \to A$，当 $\varepsilon \to 0$ 时，
即得稠密性。
\end{solution}

\begin{solution}{exo:b3:topology:9}
(a) 设 $A \subseteq \Q$ 含 $p < q$，取一个无理数
$\alpha \in (p, q)$：$A = (A \cap (-\infty, \alpha)) \sqcup
(A\cap(\alpha, +\infty))$ 把 $A$ 分成两个非空的相对
开子集，故 $A$ 不连通。连通分支皆为单点集。
局部紧性不成立：设 $V$ 为 $0$ 在 $\Q$ 中的紧邻域，
它包含某个 $[-r, r]\cap\Q$（$r > 0$），该集
在 $V$ 中为闭集，从而为紧集；但 $[-r,r]$ 中的有理数列
在 $\R$ 中收敛到一个无理数，其任何子列
都不能\emph{在 $\Q$ 中}收敛，这与
\cref{thm:b3:topology:metriccompact}矛盾。

(b) 设 $\gamma = (\gamma_1, \gamma_2) \colon [0,1] \to S$ 为
道路，满足 $\gamma(0) = (0,0)$，且 $\gamma(1) \in \Gamma$。
集合 $\{t : \gamma_1(t) = 0\}$ 为闭集，且不含
$1$；记 $t^*$ 为其上确界，于是 $\gamma_1(t^*) = 0$，且
$\gamma_1 > 0$ 于 $(t^*, 1]$ 上。由 $\gamma_2$ 在
$t^*$ 处连续，取 $\delta$，使 $\abs{\gamma_2(t) -
\gamma_2(t^*)} < \tfrac12$，对 $t \in [t^*, t^* + \delta]$ 成立。
取定 $t_0 = t^* + \delta$（不妨设 $\leq 1$）：$\gamma_1(t_0) > 0$，
且 $\gamma_1$ 取遍 $(0, \gamma_1(t_0)]$ 中的每个值，
在 $[t^*, t_0]$ 上（由介值定理，
见\cref{thm:b3:topology:connectedbasics}）。取充分大的 $k$，使
$a_k = \frac1{\pi/2 + 2k\pi} < \gamma_1(t_0)$ 且 $b_k =
\frac1{3\pi/2 + 2k\pi} < \gamma_1(t_0)$，存在 $s, s' \in
(t^*, t_0]$，使 $\gamma_1(s) = a_k$，$\gamma_1(s') = b_k$；
由于这些点位于图像上，$\gamma_2(s) = \sin(1/a_k) =
1$ 且 $\gamma_2(s') = -1$。二者不可能都位于以 $\frac12$ 为
半径、以 $\gamma_2(t^*)$ 为中心的邻域内：矛盾。$S$ 连通，却
不是道路连通的。
\end{solution}

\begin{solution}{exo:b3:topology:10}
(a) $C_n$ 恰由可写成三进制展开的 $x$ 组成，
其数字取自 $\{0, 2\}$，直到第 $n$ 位（归纳法：
去掉中间三分之一即去掉首位数字 $1$，依此类推；端点
有两种展开，其中之一不含数字 $1$），于是 $C = \bigcap C_n$
是所有和 $\sum a_n3^{-n}$（$a_n \in \{0,2\}$）构成的集合。紧性：
每个 $C_n$ 是有限个闭区间的并；$C$ 是
$[0,1]$ 中的闭集。内部为空：$C \subseteq C_n$，它是
长为 $3^{-n}$ 的区间之并。若存在长为
$\ell > 0$ 的内部区间，则对所有 $n$ 它都落在其中之一内。完全集：给定
$x \in C$ 与 $m$，翻转数字 $a_m$，所得新点仍在
$C$ 中，不同于原点，且两者距离在 $2\cdot3^{-m}$ 以内（原点即 $x$）。

(b) 若 $x \neq y \in C$，则二者的展开首次不同
发生在第 $k$ 位；二者之间有一个被去掉的中间三分之一区间
（第 $k$ 位上分隔数字 $0$ 与数字 $2$ 的间隙），
从而得到一点 $c \notin C$，不妨设 $x < c < y$，于是 $C =
(C \cap (-\infty, c)) \sqcup (C \cap (c, \infty))$ 把任一同时包含这两点的子集分开。
连通分支都是单点集。

(c) 数字映射 $\beta\colon x \mapsto (a_n(x))$ 是
到 $\{0,2\}^\N$ 上的双射（存在性与唯一性：
$\{0,2\}$ 展开中，不同的数字序列所给出的点，
在首次不同之处的距离 $\geq 3^{-k} > 0$，参见
\cref{pb:b3:topology:1} 第 1 题）。映射连续：
$\abs{x - y} < 3^{-m}$ 蕴含前 $m$ 位数字相同
（同样的计算），故 $\beta$ 把小球映入拓扑基
中的基本方块。从紧集 $C$ 到
豪斯多夫乘积空间的连续双射是同胚
（\cref{cor:b3:topology:compacthomeo}；乘积是豪斯多夫的：
在不同坐标处分离即可）。不可数性：康托的
对角线法作用于 $\{0,2\}^\N$。最后，$\{0,2\}^\N \times
\{0,2\}^\N \cong \{0,2\}^\N$（通过交错数字，这是一个
同胚：两个方向上分量均连续），故 $C
\times C \cong C$。
\end{solution}

\begin{solution}{exo:b3:topology:11}
对 $y \in \R^n$，令
\[
\tau(y) = \Bigl(\frac{2y}{\norm y^2 + 1},\ \frac{\norm y^2 -
1}{\norm y^2 + 1}\Bigr) \in \R^n\times\R = \R^{n+1}.
\]
直接验证 $\norm{\tau(y)} = 1$，$\tau(y) \neq N$，并且
$\sigma(\tau(y)) = y$，$\tau(\sigma(x)) = x$：$\sigma$ 与
$\tau$ 互逆，二者都连续（有理式，分母
不为零）：$S^n\setminus\{N\} \cong
\R^n$。将 $\hat\sigma \colon S^n \to \widehat{\R^n}$ 延拓为
$N \mapsto \infty$：它是双射，且在 $S^n\setminus\{N\}$ 的每一点处连续，
并在 $N$ 处：$\infty$ 的基本邻域为
$\widehat{\R^n}\setminus K$，$K$ 紧，$K
\subseteq \bar B(0, R)$。由于 $\norm{\sigma(x)}^2 = \frac{1 +
x_{n+1}}{1 - x_{n+1}} \to \infty$（当 $x_{n+1} \to 1$ 时），集合
$\{x \in S^n : x_{n+1} > 1 - \eta\}$ 映到 $\bar B(0,R)$ 之外，
对充分小的 $\eta$，连续性成立。从紧集
$S^n$ 到豪斯多夫空间 $\widehat{\R^n}$ 的连续双射是
同胚。去掉另一点 $p$：$S^n$ 的一个旋转把
$p$ 映到 $N$（旋转是同胚），化为已算过的情形：
$S^n \setminus \{p\} \cong \R^n$。
\end{solution}

\begin{solution}{exo:b3:topology:12}
(a) $T$ 有界且是闭集：$T$ 中点列的极限，
若其横坐标 $\to x_0 > 0$，则仍在（局部闭的）图像上，
这由 $\sin\frac1x$ 在 $\intoc01$ 上连续得到；横坐标
$\to 0$ 的极限，其纵坐标属于 $\intcc{-1}1$，因而
落在线段上。由海涅–博雷尔定理
（\cref{cor:b3:topology:heineborel}）得紧。图像
$\Gamma$ 的闭包：每一点 $(0, y)$（$y \in \intcc{-1}1$）都是
图像点的极限 --- 求解 $\sin\frac1x = y$（在 $0$ 附近）；
（该函数在每个区间上扫过 $\intcc{-1}1$，
$[\frac1{(k+1)\pi}, \frac1{k\pi}]$）：$\bar\Gamma = T$。

(b) $\Gamma$ 是连通集
$\intoc01$ 在 $x \mapsto (x, \sin\frac1x)$ 下的连续像，因而连通；
而连通集的闭包仍连通
（\cref{thm:b3:topology:connectedbasics}）：$T = \bar\Gamma$
连通。

(c) 设 $\gamma\colon\intcc01\to T$ 连续，且
$\gamma(0) = (0,0)$，$\gamma(1) = (1, \sin1)$，并令 $t^*
= \sup\{t : \gamma_1(t) = 0\}$：由连续性，$\gamma_1(t^*)
= 0$，且 $\gamma_1 > 0$ 在 $\intoc{t^*}1$ 上。对每个
$\delta > 0$，在区间 $\intoc{t^*}{t^*+\delta}$ 上
$\gamma_1$ 取遍某个 $\intoc0\eta$ 中的全部值
（由介值定理，$\gamma_1(t^* + \delta) > 0$）；
特别地，它包含形如
$\frac1{\pi/2 + 2k\pi}$ 与 $\frac1{-\pi/2 + 2k\pi}$ 的
横坐标，其中 $k$ 可任意大，此时 $\gamma_2 = \pm1$。因此
在 $t^*$ 的每个右邻域内，$\gamma_2$ 同时取到
值 $1$ 与 $-1$：$\gamma_2$ 在 $t^{*+}$ 处无极限，
与连续性矛盾。因此不存在这样的道路。

(d) $T$ 连通但不道路连通：两个概念并不相同。
对连通开集 $U \subseteq \R^n$ 与 $a \in U$，
令 $V$ 为 $U$ 中可与 $a$ 用道路相连的点的集合，
该道路位于 $U$ 内。$V$ 是开集：对 $x \in V$，存在球 $B(x, r)
\subseteq U$，它是星形的，把通向 $x$ 的道路与一条线段接起来，
便可到达球中每一点。$V$ 在
$U$ 中是闭的：若 $x \in U \setminus V$，同一球
论证表明 $B(x, r) \cap V = \varnothing$（球中若有 $V$ 的点，
则它会把 $a$ 与 $x$ 连起来）。非空（$a \in
V$），在连通的 $U$ 中既开又闭：$V = U$。正弦
曲线之所以能避开这一结论，是因为它是闭的且内部为空：其
“坏”点 $(0,0)$ 在 $T$ 中没有球可用来跨越
振荡。
\end{solution}

\begin{solution}{pb:b3:topology:1}
\textbf{1.} 设 $x, y \in C$ 的展开首次
在第 $k \leq m$ 位不同，不妨设 $b_k(x) = 1$，$b_k(y) = 0$。于是
\[
x - y \;\geq\; \frac{2}{3^k} - \sum_{j > k}\frac{2}{3^j}
= \frac{2}{3^k} - \frac{1}{3^k} = \frac1{3^k} \geq \frac1{3^m}:
\]
逆否命题：$\abs{x - y} < 3^{-m}$ 蕴含
前 $m$ 位数字一致。

\textbf{2.} 由第 1 题，$b_n$ 在 $C \cap (x -
3^{-n}, x + 3^{-n})$ 上为常值，故局部常值，从而连续。
映射 $x \mapsto (b_n(x))_n \in \{0,1\}^\N$ 是连续的，
（逐分量，由\cref{prop:b3:topology:universal}(a)），且
为双射（$\{0,2\}$ 位展开唯一）；由紧集
$C$ 到豪斯多夫空间的连续双射是同胚
（见\cref{cor:b3:topology:compacthomeo}）。

\textbf{3.} 连续性：若 $\abs{x - y} < 3^{-m}$，则前
$m$ 位数字相同，故 $\abs{g(x) - g(y)} \leq \sum_{n > m}2^{-n}
= 2^{-m}$。满射：每个 $t \in [0,1]$ 都有二进制展开
$t = \sum b_n2^{-n}$，且 $t = g\bigl(\sum
2b_n3^{-n}\bigr)$。不单射：$g$ 将两个康托点，
数字分别为 $(0,1,1,1,\dots)$ 与 $(1,0,0,\dots)$ ---
二者都映到 $\frac12$（二进制歧义 $0.0111\ldots_2 =
0.1000\ldots_2$).

\textbf{4.} 由同样的估计，$\Phi$ 的每个分量连续
（其数字为 $b_n$ 的子列）。
满射：任给 $(s, t) \in [0,1]^2$，
取 $(c_n)$ 为 $s$ 的二进制数字，$(d_n)$ 为 $t$ 的二进制数字，交错排列，点 $x
\in C$ 满足 $b_{2n-1}(x) = c_n$，$b_{2n}(x) = d_n$，则
$\Phi(x) = (s,t)$.

\textbf{5.} 各挖去的区间两两不交，端点属于
$C$，故公式唯一地确定 $f$，且定义在 $[0,1]$ 上；在挖去区间的
端点处，仿射公式给出 $\Phi(u)$，
$\Phi(v)$，与 $f = \Phi$ 在 $C$ 上相容。满射性：
已有 $f(C) = \Phi(C) = [0,1]^2$。

\textbf{6.} 当 $t_0 \notin C$ 时：$t_0$ 位于一开挖去区间内，
其上 $f$ 仿射，故连续。当 $x \in C$ 时：
先记两个估计。

\emph{(i) 在 $C$ 上}：若 $\abs{y - x} \leq 3^{-M}$，$y \in C$，
则数字直到第 $M$ 位都相同，故 $\Phi(y)
- \Phi(x)$ 的每个分量至多为 $\sum_{n > \lfloor M/2\rfloor} 2^{-n} =
2^{-\lfloor M/2\rfloor}$.

\emph{(ii) 跨越挖去区间}：区间 $(u,v)$ 于第 $k$ 步被挖去，
长为 $3^{-k}$，且其端点的数字
直到第 $k - 1$ 位都相同（自第 $k$ 位起不同），故
$\norm{\Phi(v) - \Phi(u)}_\infty \leq
2^{-\lfloor (k-1)/2\rfloor}$.

现设 $\abs{t - x} < 3^{-M}$，$t$ 位于挖去区间 $(u, v)$ 中，
其挖去于第 $k$ 步，不妨设 $u$ 为 $x$ 一侧的端点，
则 $\abs{u - x} < 3^{-M}$，且 $0 < t - u < 3^{-M}$。于是
\[
\norm{f(t) - \Phi(x)}_\infty
\leq \norm{\Phi(u) - \Phi(x)}_\infty
+ \frac{t - u}{v - u}\,\norm{\Phi(v) - \Phi(u)}_\infty
\leq 2^{-\lfloor M/2\rfloor} + \min\bigl(1,\ 3^{k - M}\bigr)\,
2^{-\lfloor(k-1)/2\rfloor} .
\]
若 $k \geq M$，则末项 $\leq 2^{-\lfloor(M-1)/2
\rfloor}$；若 $k < M$，由于 $2^{-\lfloor(k-1)/2\rfloor} \leq
\sqrt2\,\cdot 2^{-k/2}$,
\[
3^{k-M}\,2^{-\lfloor(k-1)/2\rfloor}
\leq \sqrt2\,\Bigl(\frac{3}{\sqrt2}\Bigr)^{k} 3^{-M}
\leq \sqrt2\,\Bigl(\frac{3}{\sqrt2}\Bigr)^{M} 3^{-M}
= \sqrt2\; 2^{-M/2}
\]
（中间的量关于 $k$ 递增）。总之
$\norm{f(t) - \Phi(x)}_\infty \leq C\,2^{-M/2}$
其中 $C$ 为绝对常数，令 $M \to \infty$
即得在 $x$ 处的连续性（$t \in C$ 的点由 (i) 覆盖）。因此
$f$ 是 $[0,1] \to [0,1]^2$ 的连续满射 --- 事实上
估计表明它属赫尔德型，指数相当于 $\log 2/(2\log
3)$ 量级，但这里只主张连续性。

\textbf{7.} 对 $[0,1]^n$：将 $x \in C$ 的数字
分为 $n$ 个交错的子列，并原样重复第 4--6 题的
论证。对 $\R \to \R^2$：设 $h \colon [0,1] \to [-1,1]^2$
为连续满射（将 $f$ 作伸缩即可）。定义 $F$ 于 $[m, m
+ \frac12]$（$m \in \Z$）上为 $h$ 的伸缩副本，使其充满
$[-(m+1), m+1]^2$，对 $m \geq 0$（以及对 $m <
0$ 时对称定义），在 $[m + \frac12, m + 1]$ 上取
连接两端点值的仿射线段。$F$ 连续（在闭片上粘合，
由\cref{prop:b3:topology:gluing}(b)），且其像
包含 $\bigcup_m [-(m+1), m+1]^2 = \R^2$。

\textbf{8.} $[0,1]$ 是紧的，$[0,1]^2$ 是豪斯多夫空间：
连续单射是到其像上的同胚
（对余限制应用\cref{cor:b3:topology:compacthomeo}）。

\textbf{9.} 去掉 $\frac12$ 后，$[0,1]\setminus\{\frac12\}$
不连通。若 $h \colon [0,1] \to [0,1]^2$ 为同胚，
则 $[0,1]^2\setminus\{h(\frac12)\}$ 也应
不连通（不连通空间的同胚像仍不连通）；但去掉一点的正方形
是道路连通的：用两段折线连接任意两点，
折线可避开被挖去的点。
矛盾，故 $[0,1] \not\cong [0,1]^2$。

\textbf{10.} 由第 8 题，
$[0,1] \to [0,1]^2$ 上的连续双射必为同胚，
与第 9 题矛盾。紧性恰在
\cref{cor:b3:topology:compacthomeo}中起作用：正是它使
连续双射的逆连续（闭集的像是紧的，从而是闭的）。
若无紧性，结论确实失效：$[0, 2\pi) \to S^1$ 是
连续双射，
但它不是同胚。

\textbf{11.} 本题证明：\emph{存在从 $[0,1]$ 到
$[0,1]^2$ 上的连续满射，但不存在连续双射；
特别地 $[0,1] \not\cong [0,1]^2$。}
因此“维数”在连续满射下并不保持 ---
基数乃至连续性都无法刻画它 --- 但它
\emph{确实是}拓扑不变量，在
同胚的层面上，至少对 $1$ 维与 $2$ 维而言，其中
去掉一点后的连通性便已足够。一般的陈述
--- \emph{区域不变性}：$\R^m \cong \R^n$ 蕴含 $m =
n$，且连续单射 $\R^n \to \R^n$ 是开映射 ---
此结论成立，但需要代数拓扑（同调论），超出
本课程的范围。

\textbf{12.} 在三进制下，$\frac C2 = \{\sum_na_n3^{-n} : a_n
\in \{0, 1\}\}$（将数字 $\{0,2\}$ 减半得
数字 $\{0,1\}$）。对 $t \in [0,1]$，任取三进制展开 $t =
\sum_nd_n3^{-n}$，$d_n \in \{0,1,2\}$，并将每个数字拆分
为 $d_n = a_n + b_n$，其中 $a_n, b_n \in \{0,1\}$（$0 = 0{+}0$，
$1 = 1{+}0$，$2 = 1{+}1$）：于是 $t \in \frac C2 + \frac C2$。
关键在于每个数字都\emph{在}$\{0,1\}$之内拆分，
因此不会有进位传播，各位数字可以独立处理。
反向包含显然（各位数字之和仍
$\leq 2$）。把区间伸缩 $2$ 倍，即得 $C + C = [0,2]$。

\textbf{13.} 映射 $(x, y) \mapsto x + y$ 将 $C \times C$
满射到 $C + C = [0,2]$：尘埃集不含任何正方形（内部为空，
见\cref{exo:b3:topology:10}），沿
反对角线投影到一条完整线段上。自相似性：康托点的第一个三进制数字
为 $0$ 或 $2$，去掉它后得
$C = \frac C3 \cup \bigl(\frac C3 +
\frac23\bigr)$，这正是生成 $C$ 每幅图像的
不动点方程。

\textbf{14.} 借助 $C \cong \{0,1\}^\N$（第 2 题），
交错映射 $(\varepsilon, \varepsilon') \mapsto
(\varepsilon_1, \varepsilon_1', \varepsilon_2,
\varepsilon_2', \dots)$ 是从
$\{0,1\}^\N\times\{0,1\}^\N \to \{0,1\}^\N$ 的双射，且两个方向都连续
（每个输出坐标只依赖于一个输入坐标；乘积
拓扑）。因此 $C\times C \cong C$，
并由归纳得 $C^n \cong C$。常见的空间则不然：
$\R^2 \not\cong \R$ 且 $[0,1]^2 \not\cong [0,1]$
（第 11 题）；无限乘积则具有此性质：$([0,1]^\N)^2 \cong
[0,1]^\N$，方法同样是交错映射。

\textbf{15.} $\sum_{k\geq1}2\cdot3^{-2k} = 2\cdot\frac{1/9}{1
- 1/9} = \frac14$：$\frac14$ 的三进制展开为
$0.\overline{02}$，数字属于 $\{0,2\}$，所以 $\frac14 \in
C$；它既非有限展开，又非最终为 $2$ 的尾部，
因此它不是任何被去掉的中间三分之一的端点。端点
构成可数集（每个被去掉的区间贡献两个端点，被去掉的区间可数多个），
而 $C \cong \{0,1\}^\N$ 不可数（对角线论证，或第 19 题），
因此几乎所有康托点
都像 $\frac14$ 一样，在通常的端点
图像中不可见。

\textbf{16.} $g(x) = g(y)$ 意味着两个二进制序列
$(b_n(x))$、$(b_n(y))$ 表示同一实数。同一实数的不同二进制序列
恰好出现在二进制歧义
$0.w01^\infty = 0.w10^\infty$ 处：至多两个原像，
而恰有两个原像的，正是 $(0,1)$ 中的二进有理数。
因此 $g$ 通过粘合可数多对点，将 $C$ 压缩到 $[0,1]$ 上
——这正是
\cref{ch:b3:measure}中康托阶梯的组合骨架。

\textbf{17.} $x$ 在 $P$ 中不是孤立点，故 $B(x, r/2)$
含有某点 $y_0 \in P \setminus \{x\}$；$y_0$ 也不是
孤立点，故 $B(x, r/2)$ 含有第二个点 $y_1
\in P$。任取 $\rho \leq \min\bigl(\frac{d(y_0,y_1)}3,
\frac r4\bigr)$，则 $\bar B(y_0, \rho)$ 与 $\bar B(y_1,
\rho)$ 不交，且都含于 $B(x, r)$。两者中心均在 $P$ 中，
故同样的两点提取可以重复进行：完备性
提供了取之不尽的邻近点，
使递归永远持续下去。

\textbf{18.} 对字的长度归纳定义 $F_w$：
$F_\varnothing = P$，在每个球 $B(y_w, \rho_w)$ 内，
由第 17 题得两个不交的闭球，半径 $\leq
\min(\rho_w/2,\ 2^{-\abs w-1}\operatorname{diam}P)$，中心在
$P$ 中的点处；令 $F_{wi} = \bar B(y_{wi}, \rho_{wi})
\cap P$：非空（含其中心）且紧。对无限
序列 $\varepsilon$，嵌套的非空紧集
$F_{\varepsilon_1\cdots\varepsilon_m}$ 的交非空
（由紧集 $P$ 中的有限交性质），直径为 $0$，
即单点 $h(\varepsilon)$。

\textbf{19.} 在第 $m$ 位首次不同的两个字，其像
分别落入两个\emph{不交}的集合 $F_{w0}, F_{w1}$
（公共前缀为 $w$），故 $h$ 是单射。若两个字在前
$m$ 位一致，则两个像都落在直径 $\leq
2^{-m}\operatorname{diam}P$ 的集合中，因此 $h$ 连续。于是
不可数集 $\{0,1\}^\N$ 单射入 $P$：每个完备紧度量空间
都是不可数的——$[0,1]$ 与 $C$
皆属其中。

\textbf{20.} $h$ 是从紧集 $\{0,1\}^\N$ 到豪斯多夫度量空间 $P$ 的连续单射，
由\cref{cor:b3:topology:compacthomeo}可知它是到其像上的同胚；
而 $\{0,1\}^\N \cong C$（第 2 题）。
于是每个完备紧度量空间都含有康托集的一个副本：
完备性具有一个普适的胚芽，
它便是
康托的。

\textbf{21.} 可数紧度量空间不可能是完备集（第 19 题），
故必有孤立点。在 $K =
\{0\}\cup\{\frac1n : n \geq 1\}$ 中，每个 $\frac1n$ 都是孤立点，
而 $0$ 不是：孤立点甚至可以稠密，而空间本身不离散——
紧性保证它们的极限
仍在空间之内。

\textbf{22.} 设 $P$ 为 $F$ 的凝聚点集，
$\mathcal B$ 为由有理端点开区间构成的可数族。
$F$ 的每个非凝聚点都落在
某个 $I \in \mathcal B$ 中，且 $I \cap F$ 可数，故 $F
\setminus P$ 含于这些可数个
可数迹之并中，故可数。由于 $F$ 不可数，
$P \neq \varnothing$；$P$ 为闭集（若 $x \notin P$，则
见证区间 $I$ 中不含任何凝聚点，
即 $I \cap F$ 可数）。且 $P \subseteq F$（$F$ 闭：
凝聚点特别地是附着点）。$P$ 为
完备集：若 $I \cap P = \{x\}$ 对某个区间 $I$ 成立，则
$I \cap F \subseteq \{x\} \cup (F \setminus P)$ 则为
可数的，与 $x \in P$ 矛盾。现在
第 17--19 题的树形构造原样适用于 $P$ 中：各片段
$\bar B(y, \rho) \cap P$ 皆为紧的（即 $\R$ 中的
有界闭子集），每个中心在 $P$ 中都非孤立，论证
未用到更多条件，因此 $\abs P \geq
\abs{\{0,1\}^\N} = \mathfrak c$，而 $\abs F \leq \mathfrak
c$ 显然成立：不可数闭集 $F \subseteq \R$ 的
基数恰为 $\mathfrak c$。闭集
无法见证连续统假设的失效。

\textbf{23.} 设 $x \neq y$ 为 $C$ 中的两点，取 $m \geq 0$，使
$3^{-(m+1)} \leq \abs{x - y} < 3^{-m}$。由第 1 题，
数字 $b_n(x) = b_n(y)$ 在 $n \leq m$ 时相同；
$\Phi$ 的第一个坐标用 $b_1, b_3, \dots$，第二个坐标用 $b_2,
b_4, \dots$，因此每个坐标中两点至少
共享前 $\lfloor m/2\rfloor$ 个二进制数字：
\[
\norm{\Phi(x) - \Phi(y)}_\infty
\leq \sum_{k > \lfloor m/2\rfloor}2^{-k}
= 2^{-\lfloor m/2\rfloor} \leq \sqrt2\cdot2^{-m/2}
= \sqrt2\,\bigl(3^{-m}\bigr)^\alpha
\leq \sqrt2\,\bigl(3\abs{x-y}\bigr)^\alpha,
\]
利用 $2^{-m/2} = 3^{-m\alpha}$（取对数即得）；并
$\sqrt2\cdot3^\alpha = \sqrt2\cdot\sqrt2 = 2$，于是
$2\abs{x-y}^\alpha$ 即为 $C$ 上的界。在同一余区间 $(u, v)$ 上 $f$ 为仿射函数，
因此对 $t, s \in [u, v]$，
\[
\norm{f(t) - f(s)} = \frac{\abs{t-s}}{v-u}\,
\norm{\Phi(v) - \Phi(u)}
\leq \frac{\abs{t-s}}{v-u}\,2(v-u)^\alpha
\leq 2\abs{t-s}^\alpha,
\]
由于 $\frac{\abs{t-s}}{v-u} \leq \bigl(\frac{\abs{t-s}}
{v-u}\bigr)^\alpha$ 当 $\abs{t-s} \leq v - u$ 时成立。一般地，设 $t <
s$，若 $[t, s]$ 与 $C$ 相交，记 $u'$、$v'$ 分别为
紧集 $C \cap [t, s]$ 的最小点与最大点；则 $t$
位于某个余区间内（或位于 $C$ 的某点处），该余区间的右端点为
$u'$，$s$ 位于左端点为 $v'$ 的余区间中，且
\[
\norm{f(t) - f(s)} \leq 2\abs{t - u'}^\alpha
+ 2\abs{v' - u'}^\alpha + 2\abs{s - v'}^\alpha
\leq 6\abs{t - s}^\alpha ;
\]
若 $[t, s]$ 与 $C$ 不相交，则同一余区间的情形适用。因此 $f$ 为
$\alpha$-赫尔德连续，其中 $\alpha = \frac{\ln2}{2\ln3}$。

\textbf{24.} 假设 $\norm{F(t) - F(s)} \leq
C\abs{t-s}^\alpha$，$F$ 为满射，且 $\alpha >
\frac12$。将 $[0,1]$ 分成 $N$ 个区间 $I_1, \dots, I_N$，
各长为 $\frac1N$：每个像 $F(I_j)$ 的直径 $\leq
CN^{-\alpha}$。网格 $G_M =
\bigl\{\bigl(\frac iM, \frac jM\bigr)\bigr\}_{0 \leq i, j
\leq M}$ 中任两个不同点的距离 $\geq \frac1M$，因而直径
$< \frac1M$ 的集合至多含其中一个点。取 $M =
\lfloor N^\alpha/(2C)\rfloor$（$N$ 充分大），则
$CN^{-\alpha} \leq \frac1{2M} < \frac1M$，且满射性
使 $(M+1)^2$ 个网格点中的每一个都落入某个
$F(I_j)$，由此
\[
\frac{N^{2\alpha}}{4C^2} \leq (M + 1)^2 \leq N,
\]
而 $N^{2\alpha - 1} \leq 4C^2$ 对充分大的 $N$ 不成立，
当 $2\alpha > 1$ 时即得矛盾。本例曲线的 $\alpha \approx
0.3155$，位于此界之下，这是必然的；希尔伯特曲线
（在本阶段不加证明地承认）是 $\frac12$-赫尔德连续的，因此临界指数
$\frac12$ 恰好取到 --- 赫尔德正则性不同于
单射性，它只是程度问题，而 $\frac12$ 恰是
维数 $2$ 对由
维数 $1$ 出发的映射所设定的分界。

\textbf{25.} 第 $m$ 个簇 $\{\frac1m + \frac1n : n >
m^2\}$ 位于 $\bigl(\frac1m, \frac1m + \frac1{m^2}\bigr]
\subseteq \bigl(\frac1m, \frac1{m-1}\bigr)$，当 $m \geq 2$ 时，
因为 $\frac1{m-1} - \frac1m = \frac1{m(m-1)} >
\frac1{m^2}$（且第一个簇位于 $(1, \frac32]$）：
各簇占据互不相交的区间。$K$ 的极限点：
在第 $m$ 个区间内，只有 $\frac1m$（该簇
收敛于它，且自身离散）；整体上，$0$
（$0$ 的每个邻域都包含整个簇）。因此 $K'
= \{0\} \cup \{\frac1m\}$，则 $K'' = \{0\}$（每个 $\frac1m$
在 $K'$ 中是孤立点），$K''' = \varnothing$；$K$ 包含其
全部极限点，因而是闭集、有界、可数且紧，
而其孤立点 --- 即簇点 $\frac1m +
\frac1n$ --- 在 $K$ 中稠密：$K'$ 中每个元素都是
它们的极限，正如第 21 题的机制所预言。归纳：
$K_1 = \{0\}\cup\{\frac1n\}$，有 $K_1^{(1)} = \{0\}$；
设可数紧集 $K_j \subseteq [0, 1]$，满足 $K_j^{(j)} =
\{0\}$ 且 $K_j^{(j+1)} = \varnothing$，令
\[
K_{j+1} = \{0\} \cup \bigcup_{m \geq 1}
\Bigl(\frac1m + \varepsilon_m K_j\Bigr),
\]
取 $\varepsilon_m$ 充分小，使第 $m$ 个副本位于
$\bigl(\frac1m, \frac1{m-1}\bigr)$ 内。导集运算逐副本作用
（各副本位于互不相交的开区间中），故
$K_{j+1}^{(i)} = \{0\} \cup \bigcup_m\bigl(\frac1m +
\varepsilon_mK_j^{(i)}\bigr)$，对 $i \leq j$ 成立；当 $i = j$ 时，此式
化为 $\{0\} \cup \{\frac1m\}$，则 $K_{j+1}^{(j+1)} =
\{0\}$，且 $K_{j+1}^{(j+2)} = \varnothing$。每个有限阶
都会出现。完备集则是另一极端：$P' = P$，导集
运算永不移动 --- 而康托–本迪克松（第
22 题）恰好说明：每个闭集都分解为
一个完备核心（对导集运算不可见），以及一个
可数的剩余部分，它被导集运算逐步消去。
\end{solution}
